\section{Introduction}
\label{sec:introduction}

\subsection{Motivation}

Modern data analytics increasingly relies on data lakes that integrate
multiple heterogeneous sources rather than on a single, centrally designed
relational database. Although this architecture simplifies the storage and
integration of large and diverse datasets, it also makes analytical access
more complex. A user who wants to formulate a cross-source query must know
which datasets are relevant, understand their schemas and semantics, identify
the relationships between them, and finally express the analysis using the
query language supported by the processing engine.

Large language models (LLMs) provide a promising interface for reducing this
complexity by translating natural-language analytical requests into SQL.
However, applying Text-to-SQL techniques directly to a heterogeneous data lake
introduces an additional problem: the model needs sufficient metadata to
identify the correct physical datasets, columns, relationships, and semantic
rules. Providing the complete metadata catalog for every query is simple, but
it can introduce large amounts of irrelevant information and causes the prompt
to grow as the catalog grows. Moreover, a metadata catalog may contain
semantic concepts, aliases, relationship identifiers, and descriptive rules
that are useful for retrieval but are not valid physical SQL identifiers.

This project investigates an alternative approach in which the metadata
provided to the LLM is selected dynamically for each user question. The goal
is to preserve the information required for correct query generation while
avoiding the cost and ambiguity of exposing the complete catalog.


\subsection{Problem Statement}

The project considers a heterogeneous data lake containing New York City
urban-mobility and weather data. The core sources are NYC TLC Yellow Taxi trip
records, NYC TLC Green Taxi trip records, NYC Taxi Zones, and hourly weather
observations. These datasets differ in schema, granularity, semantics, and
analytical role, while meaningful user requests may require relationships
across multiple sources.

For example, a request concerning the borough with the largest number of taxi
pickups during rainy hours requires more than simply identifying columns whose
names are similar to words in the question. The system must identify the taxi
trip dataset, the pickup-location relationship with the Taxi Zones dataset,
the temporal relationship with hourly weather observations, and the semantic
condition used to represent a rainy hour. The generated query must then use
the corresponding physical Spark SQL tables and columns rather than internal
metadata identifiers.

The central challenge is therefore to determine how much metadata should be
provided to the language model and how that metadata should be represented.
Too little information may cause the model to omit a required dataset,
relationship, or semantic rule. Too much information increases prompt size and
may expose irrelevant or misleading metadata.


\subsection{Proposed Approach}

The implemented system follows a metadata-aware Text-to-Spark-SQL pipeline.
A structured metadata catalog describes datasets, physical columns, semantic
concepts, aliases, cross-source relationships, and reusable semantic rules.
Metadata documents are indexed using vector embeddings so that candidate
metadata can be retrieved from a natural-language question.

The initial semantic retrieval is then expanded using explicit relationships
stored in the catalog. This relation-aware expansion makes it possible to add
the physical endpoints required for joins even when they are not directly
retrieved by semantic similarity. The selected metadata is subsequently
rendered into a SQL-safe context that exposes only executable physical table
and column names, join conditions, and SQL expressions to the language model.

A local LLM generates Spark SQL from this restricted context. The resulting
query is parsed and validated before execution. If validation fails, the
system can perform a single repair attempt using the validation error together
with the same retrieved metadata context. Valid queries are finally executed
using Apache Spark SQL over the curated data lake.


\subsection{Research Question}

The main research question investigated in this work is:

\begin{quote}
\emph{To what extent can dynamic, relation-aware metadata retrieval improve
the correctness and efficiency of Text-to-Spark-SQL generation over a
heterogeneous data lake compared with providing the complete metadata catalog
to the language model?}
\end{quote}

The evaluation considers both correctness and efficiency. Correctness is
measured at the metadata-retrieval and query-result levels, while efficiency is
evaluated through prompt size, controlled inference latency, catalog growth,
and Spark execution under increasing data volume.


\subsection{Contributions}

The main contributions of the project are the following:

\begin{itemize}

    \item A reproducible metadata-aware Text-to-Spark-SQL prototype for
    cross-source analytics over a heterogeneous Spark-based data lake.

    \item A structured metadata catalog containing physical schemas, semantic
    concepts, aliases, cross-source relationships, and SQL-oriented semantic
    rules.

    \item A relation-aware retrieval strategy that combines semantic retrieval
    with deterministic metadata expansion and produces a restricted,
    SQL-safe context for the language model.

    \item A validation and one-shot repair stage that prevents invalid SQL
    from being executed directly and attempts to recover from generation
    errors without changing the retrieved metadata.

    \item A frozen held-out evaluation comparing three configurations:
    complete-catalog prompting, relation-aware metadata retrieval, and
    relation-aware retrieval with validation and repair.

    \item Separate scalability experiments for metadata-catalog growth and
    physical data volume, allowing catalog scalability and Spark execution
    scalability to be analysed independently.

\end{itemize}


\subsection{Summary of Results}

On the frozen held-out benchmark, the complete-catalog baseline achieved
70\% result accuracy. Relation-aware metadata retrieval increased result
accuracy to 80\%, while the addition of validation and one-shot repair reached
85\%. At the same time, relation-aware prompting reduced the mean prompt size
from 2558.45 to 776.35 tokens, corresponding to a 69.66\% reduction.

A separate controlled latency experiment showed that the smaller
relation-aware prompt reduced mean prompt-evaluation time from 182.58 seconds
to 54.24 seconds in the evaluated CPU-only local environment, a reduction of
70.29\%.

The catalog-scaling experiment showed that the full metadata context increased
from 9,498 characters with four sources to 16,651 characters with sixteen
sources, whereas the mean retrieved context remained close to 2,100
characters. This increased compression was accompanied by a measurable
retrieval trade-off: macro recall decreased from 0.925 to 0.825 after
semantically related hard-negative sources were introduced and then remained
stable when the catalog was further expanded.

Finally, Spark data-volume experiments evaluated the same workloads from
100,000 rows up to the full Yellow Taxi dataset of 2,963,713 curated records.
These experiments complement the catalog-scale evaluation by measuring the
system along the physical data-volume dimension rather than the metadata
dimension.

The remainder of this report describes the data sources, metadata model,
system architecture, implementation, experimental methodology, results, and
limitations in detail.
